%==============================================================
%  ou.tex
%  hasklib — mathematical companion / derivations
%
%  Ornstein-Uhlenbeck: exact Gaussian solution via the integrating-
%  factor trick (It\^o's lemma applied to e^{kt}X_t), mean/variance,
%  and the stationary law.
%
%  Depends on ito_foundations.tex (It\^o's lemma, Theorem 1).
%  Sibling of gbm.tex and cir.tex, not a dependent of either.
%==============================================================
\documentclass[11pt]{article}

%---------------- packages ----------------
\usepackage[utf8]{inputenc}
\usepackage[T1]{fontenc}
\usepackage{amsmath, amssymb, amsthm}
\usepackage{mathtools}
\usepackage[a4paper, margin=1in]{geometry}
\usepackage{xcolor}
\usepackage{hyperref}
\usepackage{cleveref}
\usepackage{bm}
\usepackage{harvard}   % Harvard author-date referencing (with agsm.bst)

%---------------- notation macros ----------------
% (kept identical to ito_foundations.tex / gbm.tex so the files read as one document)
\newcommand{\Prob}{\mathbb{P}}
\newcommand{\Qmeas}{\mathbb{Q}}
\newcommand{\E}{\mathbb{E}}
\newcommand{\Var}{\operatorname{Var}}
\newcommand{\Cov}{\operatorname{Cov}}
\newcommand{\filt}{\mathcal{F}}
\newcommand{\Om}{\Omega}

\newcommand{\dd}{\mathrm{d}}
\newcommand{\dW}{\dd W_t}
\newcommand{\dt}{\dd t}
\newcommand{\R}{\mathbb{R}}

\newcommand{\drift}{a}
\newcommand{\diff}{b}

%---------------- theorem environments ----------------
\theoremstyle{plain}
\newtheorem{theorem}{Theorem}
\newtheorem{lemma}{Lemma}
\newtheorem{proposition}{Proposition}
\newtheorem{corollary}{Corollary}

\theoremstyle{definition}
\newtheorem{definition}{Definition}

\theoremstyle{remark}
\newtheorem{remark}{Remark}

\newcommand{\TODO}[1]{\medskip\noindent\textbf{[\textcolor{red}{TO DERIVE}: #1]}\medskip}

%==============================================================
\title{The Ornstein--Uhlenbeck Process: Exact Solution and Moments\\
       \large \texttt{hasklib} mathematical companion}
\author{Hubert}
\date{\today}
%==============================================================

\begin{document}
\maketitle

\begin{abstract}
This note derives the result \texttt{OU::sample\_terminal} relies on:
the Ornstein--Uhlenbeck (OU) SDE is linear enough that It\^o's lemma
applied to the integrating factor $e^{\kappa t}X_t$ removes the
$X_t$-dependence entirely, giving an exact Gaussian terminal law rather
than requiring discretisation. We derive that law, its mean and
variance, and the stationary distribution as $t\to\infty$. This file
assumes It\^o's lemma as established in \texttt{ito\_foundations.tex}
(Theorem~1 there) and a standard fact about Wiener integrals (the
It\^o isometry, cited below); neither is re-derived here.
\end{abstract}

%--------------------------------------------------------------
\section{The OU SDE}
%--------------------------------------------------------------

\begin{definition}[Ornstein--Uhlenbeck process]\label{def:ou}
$X_t$ follows an Ornstein--Uhlenbeck process with mean-reversion speed
$\kappa>0$, long-run mean $\theta\in\R$, and volatility $\sigma>0$ if
it solves
\begin{equation}\label{eq:ousde}
   \dd X_t = \kappa(\theta - X_t)\,\dt + \sigma\,\dW, \qquad X_0\in\R.
\end{equation}
In the notation of \texttt{ito\_foundations.tex} this is the choice
$\drift(t,x)=\kappa(\theta-x)$, $\diff(t,x)=\sigma$ (constant) in the
general SDE $\dd X_t = \drift(t,X_t)\dt+\diff(t,X_t)\dW$.
\end{definition}

\begin{remark}[Correspondence with the code]
\texttt{OU::drift} returns $\kappa(\theta-x)$ and
\texttt{OU::diffusion} returns the constant $\sigma$, independent of
$x$ --- unlike GBM, the diffusion here does not scale with the state.
\texttt{EulerMaruyama} can still step this process through the same
\texttt{StochasticProcess} interface, but as with GBM there is an
exact alternative, derived below, that \texttt{OU::sample\_terminal}
uses instead.
\end{remark}

Both coefficients are affine (in particular globally Lipschitz and of
linear growth) in $x$, so the existence--uniqueness assumption of
\texttt{ito\_foundations.tex} applies unconditionally here --- unlike
GBM, there is no positivity constraint to track, since OU is free to
take any real value.

%--------------------------------------------------------------
\section{Exact solution: the integrating-factor trick}
%--------------------------------------------------------------

Unlike GBM, taking a logarithm does nothing useful here since
\eqref{eq:ousde} is already linear in $X_t$, not multiplicative. The
relevant substitution instead is the same one used for linear ODEs:
multiply by the integrating factor $e^{\kappa t}$ and let It\^o's
lemma absorb the product rule.

\begin{theorem}[Exact terminal law of OU]\label{thm:ou-exact}
Let $X_t$ solve \eqref{eq:ousde}. Then for any $T>0$,
\begin{equation}\label{eq:ouexact}
   X_T = X_0\,e^{-\kappa T} + \theta\big(1-e^{-\kappa T}\big)
         + \sigma\, e^{-\kappa T}\!\int_0^T e^{\kappa s}\,\dd W_s .
\end{equation}
Conditional on $X_0$, $X_T$ is Gaussian.
\end{theorem}

\begin{proof}
Apply It\^o's lemma (Theorem~1 of \texttt{ito\_foundations.tex}) to
$f(t,x)=e^{\kappa t}x$, so $f_t=\kappa e^{\kappa t}x$, $f_x=e^{\kappa
t}$, $f_{xx}=0$. With $\drift(t,X_t)=\kappa(\theta-X_t)$ and
$\diff(t,X_t)=\sigma$,
\[
   \dd\big(e^{\kappa t}X_t\big)
     = \Big(\kappa e^{\kappa t}X_t + \kappa(\theta-X_t)\,e^{\kappa t}
             + 0\Big)\dt + \sigma\,e^{\kappa t}\,\dW .
\]
The $X_t$ terms cancel completely --- $\kappa e^{\kappa t}X_t$ against
$-\kappa X_t\,e^{\kappa t}$ --- which is precisely the point of the
integrating factor, leaving
\begin{equation}\label{eq:oulinear}
   \dd\big(e^{\kappa t}X_t\big) = \kappa\theta\,e^{\kappa t}\,\dt
                                   + \sigma\,e^{\kappa t}\,\dW .
\end{equation}
The right-hand side no longer involves $X_t$ at all, so
\eqref{eq:oulinear} integrates directly from $0$ to $T$:
\[
   e^{\kappa T}X_T - X_0
     = \kappa\theta\int_0^T e^{\kappa s}\,\dd s
       + \sigma\int_0^T e^{\kappa s}\,\dd W_s
     = \theta\big(e^{\kappa T}-1\big)
       + \sigma\int_0^T e^{\kappa s}\,\dd W_s .
\]
Multiplying by $e^{-\kappa T}$ gives \eqref{eq:ouexact}. The stochastic
integral $\int_0^T e^{\kappa s}\,\dd W_s$ is a Wiener integral of a
deterministic integrand, hence Gaussian (a linear functional of a
Gaussian process); $X_T$ is an affine transformation of it and of the
deterministic terms, so $X_T$ is Gaussian as well.
\end{proof}

\begin{lemma}[Moments of a Wiener integral]\label{lem:ito-isometry}
For deterministic $g\in L^2[0,T]$, the stochastic integral
$I=\int_0^T g(s)\,\dd W_s$ satisfies $\E[I]=0$ and
$\Var(I)=\int_0^T g(s)^2\,\dd s$ (the It\^o isometry).
\end{lemma}

\begin{proof}[Citation]
This is a standard property of the It\^o integral, following from its
construction as an $L^2$-limit of Riemann-type sums of independent
Gaussian increments; see \citeasnoun{shreve2004} or
\citeasnoun{oksendal2003}. \emph{(Cited, not proved, in the same spirit
as the existence--uniqueness assumption of \texttt{ito\_foundations.tex}
--- the library uses the conclusion, not the construction.)}
\end{proof}

%--------------------------------------------------------------
\section{Moments and the stationary law}
%--------------------------------------------------------------

\begin{proposition}[Mean and variance of OU]\label{prop:ou-moments}
\[
   \E[X_T\mid X_0] = X_0\,e^{-\kappa T} + \theta\big(1-e^{-\kappa T}\big),
   \qquad
   \Var(X_T\mid X_0) = \frac{\sigma^2}{2\kappa}\Big(1-e^{-2\kappa T}\Big).
\]
\end{proposition}

\begin{proof}
Take expectations in \eqref{eq:ouexact}. By \cref{lem:ito-isometry}
with $g(s)=\sigma e^{-\kappa T}e^{\kappa s}$ (constant $e^{-\kappa T}$
pulled outside the integral), $\E\big[\int_0^T e^{\kappa
s}\dd W_s\big]=0$, so the mean is just the deterministic part:
\[
   \E[X_T\mid X_0] = X_0\,e^{-\kappa T} + \theta\big(1-e^{-\kappa T}\big).
\]
For the variance, only the stochastic term contributes (the rest of
\eqref{eq:ouexact} is deterministic given $X_0$):
\[
   \Var(X_T\mid X_0)
     = \sigma^2 e^{-2\kappa T}\,\Var\!\left(\int_0^T e^{\kappa s}\,\dd W_s\right)
     = \sigma^2 e^{-2\kappa T}\int_0^T e^{2\kappa s}\,\dd s ,
\]
using \cref{lem:ito-isometry} with $g(s)=e^{\kappa s}$. The integral
evaluates to $\big(e^{2\kappa T}-1\big)/(2\kappa)$, giving
\[
   \Var(X_T\mid X_0)
     = \sigma^2 e^{-2\kappa T}\cdot\frac{e^{2\kappa T}-1}{2\kappa}
     = \frac{\sigma^2}{2\kappa}\Big(1-e^{-2\kappa T}\Big). \qedhere
\]
\end{proof}

\begin{remark}
Both the mean and the variance are exactly what
\texttt{OU::sample\_terminal} needs: draw $Z\sim\mathcal N(0,1)$ once
and set
\[
   X_T = X_0\,e^{-\kappa T} + \theta\big(1-e^{-\kappa T}\big)
         + \sqrt{\frac{\sigma^2}{2\kappa}\big(1-e^{-2\kappa T}\big)}\;Z,
\]
exactly as GBM's exact sampler needed only $Z$ and the closed-form mean
and variance from \texttt{gbm.tex}. No time-stepping, hence no
discretisation bias.
\end{remark}

\begin{corollary}[Stationary law]\label{cor:ou-stationary}
As $T\to\infty$ (with $\kappa>0$ fixed), $X_T\mid X_0$ converges in
distribution to
\[
   X_\infty \sim \mathcal N\!\left(\theta,\ \frac{\sigma^2}{2\kappa}\right),
\]
independent of the starting point $X_0$.
\end{corollary}

\begin{proof}
Immediate from \cref{prop:ou-moments}: $e^{-\kappa T}\to 0$ and
$e^{-2\kappa T}\to 0$ as $T\to\infty$ since $\kappa>0$, so the mean
$\to\theta$ and the variance $\to\sigma^2/(2\kappa)$, independent of
$X_0$.
\end{proof}

\begin{remark}[Why $\kappa>0$ matters]
The mean-reversion speed $\kappa>0$ is what makes $e^{-\kappa T}\to0$
and hence guarantees a stationary law at all; it is also what makes OU
a sensible model for a quantity that fluctuates around a long-run
level $\theta$ rather than drifting away from it, in contrast to GBM's
persistent (possibly explosive) drift. If $\kappa\le0$ the process is
non-stationary and \cref{cor:ou-stationary} does not apply --- the
library assumes $\kappa>0$ throughout.
\end{remark}

\begin{remark}[Relevance beyond the pricing showcase]
OU is not only one of the four processes in the convergence
demonstration; it is also the driving process behind the Hull--White
short-rate model, $\dd r_t = \kappa(\theta(t)-r_t)\dt+\sigma\,\dW$,
used throughout \citeasnoun{brigomercurio2006} for term-structure
modelling. The only difference from \eqref{eq:ousde} is a
time-dependent long-run level $\theta(t)$ calibrated to the observed
yield curve, chosen precisely so that the affine structure exploited in
\cref{thm:ou-exact} survives with $\theta$ replaced by $\theta(t)$
inside the integral. This is the model relevant to the callable-note
decomposition (cancelable payer swap plus receiver swaption) discussed
elsewhere in this thesis.
\end{remark}

%--------------------------------------------------------------
\bibliographystyle{agsm}
\bibliography{refs}

\end{document}